\section{Related work}
\label{sec:related}
\paragraph{Aggregation and evaluation.} Multi-task benchmarks such as MMLU~\citep{hendrycks2021mmlu}, BIG-Bench Hard~\citep{suzgun2022bbh} and the HELM framework~\citep{liang2022helm} report a per-task score matrix that is then aggregated, usually by a (weighted) mean. \citet{dehghani2021lottery} show that task selection and aggregation choices can change the relative ordering of methods (``benchmark lottery''), and \citet{lanctot2023socialchoice} propose voting-theoretic aggregation for multi-task evaluation of agents and language models. These works treat the rule as the evaluator's choice; we instead ask what a \emph{developer} must spend to win under each rule.

\paragraph{Strategic behaviour on leaderboards and contamination.} \citet{singh2025illusion} analyse private testing and sampling asymmetries on Chatbot Arena; \citet{zhou2023cheater} and \citet{yang2023rephrased} study training on benchmark data and the failure of decontamination, the real-world counterpart of ``targeted investment''. \citet{chen2026incentives} model benchmarking as a Stackelberg game between a designer and several competing developers who allocate post-training resources and characterise equilibria. We study a different, complementary object: the minimum cost for \emph{one} target to be first, rule by rule, measured empirically on public leaderboards, with no equilibrium reasoning.

\paragraph{Computational social choice.} Making a designated candidate win by changing the input at minimum cost is the bribery problem~\citep{faliszewski2009bribery,faliszewski2009llull,faliszewski2016control,brandt2016handbook}, and shift bribery and margin of victory are close relatives. Here tasks play the role of voters, but the manipulated object is the candidate's own \emph{cardinal} score vector (not voters' preferences), the number of voters $n$ is small, costs are continuous and capped, and the question is empirical on real score matrices.
